\subsection{Core Set Mutation Operations}

Set mutators (\texttt{add}, \texttt{remove}, \texttt{discard}, \texttt{pop}, \texttt{clear}) update the same \texttt{PySetObject} in place; method return values are \texttt{None} except \texttt{pop}, which returns the removed element.

\subsubsection{Adding Elements with \texttt{.add()}}

\texttt{add} hashes the item, probes, and inserts only if no equal element exists. Duplicate adds are silent no-ops. Unhashable objects fail at hash time with \texttt{TypeError}.

\subsubsection{Removing Elements with \texttt{.remove()}, \texttt{.discard()}, \texttt{.pop()}, and \texttt{.clear()}}

\texttt{remove} raises \texttt{KeyError} on miss; \texttt{discard} ignores misses. \texttt{pop} removes an arbitrary element---sets expose no first/last ordering. Deletion may leave dummy probe markers. \texttt{clear} resets membership while preserving object identity.

\subsubsection{Membership Testing with \texttt{in}}

\texttt{in} executes the standard hash-then-\texttt{\_\_eq\_\_} probe sequence. Both the set and the searched value must be hashable.
